\section{Finite local weights for retained rows}
\label{sec:retained-rows-v42}

\status{This chapter constructs the original finite-group idempotents,
proves positivity and normalization of their original series weights,
and proves the squarefree, divisor-degree, long-prime-degree, and
root-to-twist collision parts of the local probability estimate.
It also constructs the complete
row labeling from exact additive dimensions and the prescribed good
rows.  The squarefree character identity, the comparison of the two
local characters, the root-injectivity collision bound, and the actual
automorphic row/count identifications remain supplied inputs.  No
whole global obligation is closed.}

All source references are to
\texttt{7cbf7a0c32fca1badf05931879a37c3f8a2b03b4}.
The index $k\in\{0,1\}$ used for the two auxiliary primes is kept
separate from the integer $k_0$ in the finite unipotent subgroup.
The group-theoretic arguments below work over any finite field
$\mathbb F_q$; in the original application $q=g_k$ is prime,
$N=R.r$, $k_0=R.k$, and $s=R.m/2-1$.  The original rank conditions
give $N=k_0+2s+1$ and $k_0\ge1$.

\subsection{The actual finite idempotents and their probability laws}

Put $G=\operatorname{GL}_N(\mathbb F_q)$ and let $U$ be its upper
unitriangular subgroup.  All matrix indices start at zero.  Let
$W=\{s,\ldots,s+k_0-1\}$, $c=s+k_0$, and
\[
 J=\{u\in U:u_{a,c}=0\quad(a\in W)\}.
 \tag{RRT.1}
\]
For a nontrivial additive character $\psi:\mathbb F_q\to\mathbb C^\times$,
let $\phi_U(u)=\sum_{i=0}^{N-2}u_{i,i+1}$.
The function $\phi_J$ sums the entries along each of two ordered
chains: the chain $s,s+1,\ldots,s+k_0-1$, and the chain
\[
                  0,1,\ldots,s-1,c,c+1,\ldots,N-1.
\]
When $s=0$, the second chain is the singleton $c$ and contributes
zero.  This is exactly the source's \texttt{phiJ}, including its
truncated sums when $s=0$ or $k_0=1$.  Define in $\mathbb C[G]$
\[
 e_U=|U|^{-1}\sum_{u\in U}\psi(\phi_U(u))^{-1}[u],\qquad
 e_J=|J|^{-1}\sum_{u\in J}\psi(\phi_J(u))^{-1}[u].
 \tag{RRT.2}
\]
These are the literal \texttt{eU} and \texttt{eJ}, not conjugate
surrogates.  Write $[1]a$ for the coefficient of the identity in
$a\in\mathbb C[G]$.

% BEGIN RRT1
\paragraph{RRT1. Exact idempotents and identity coefficients.}
\label{claim:RRT1}
Writing $M=N(N-1)/2$, we have
\[
 e_U^2=e_U,\quad e_J^2=e_J,\qquad
 [1]e_U=q^{-M},\quad [1]e_J=q^{k_0-M}>0.
 \tag{RRT.3}
\]
Both idempotents are self-adjoint for the involution
$a^*(g)=\overline{a(g^{-1})}$.

Indeed, for upper unitriangular $u,v$ and $a<b$,
\[
 (uv)_{a,b}=u_{a,b}+v_{a,b}
                  +\sum_{a<h<b}u_{a,h}v_{h,b}.
 \tag{RRT.4}
\]
For $a\in W,b=c$, every summand in the last sum is zero since
$h\in W$ and $v_{h,c}=0$.  Thus $J$ contains the identity and is
closed under multiplication; as a finite subset of a group, it is a
subgroup.  All entries in the two character sums are adjacent,
except possibly $(s-1,c)$.  For that entry the intervening indices
are precisely $W$, so the same calculation makes its cross terms
zero.  Consequently $\phi_U$ and $\phi_J$ are additive under
group multiplication on their respective subgroups.

For any finite subgroup $H$ and character $\lambda:H\to\mathbb C^\times$,
write $e=|H|^{-1}\sum_h\lambda(h)^{-1}[h]$.
The coefficient of $[z]$ in $e^2$ is zero outside
$H$, and for $z\in H$ it is
$|H|^{-2}\sum_{h\in H}\lambda(z)^{-1}=|H|^{-1}\lambda(z)^{-1}$.
This proves idempotence.  A finite-order complex character has
absolute value one, which proves self-adjointness.  Its coefficient
at $1$ is $|H|^{-1}$.  Finally $U$ has $M$ freely chosen entries,
and $J$ has exactly $M-k_0$ freely chosen entries: the $k_0$
forbidden entries in (RRT.1) are distinct and lie above the diagonal.
This gives (RRT.3).  No Green parametrization is used.
\hfill$\square$
% END RRT1

% BEGIN RRT2
\paragraph{RRT2. The original series weights are probabilities.}
\label{claim:RRT2}
Let $E_F$ be the actual central series idempotents
$D.\mathrm{ntypes}.\mathrm{eSer}(k,F)$ for a fixed auxiliary prime.
They are mutually orthogonal, vanish unless $F$ is monic of degree
$N$ with nonzero constant coefficient, and sum to $1$ over those
labels.  These properties are explicit fields of the original
\texttt{NTypes}.  For $a=e_U$ or $e_J$, set
\[
             w_a(F)=\frac{\operatorname{Re}[1](aE_F)}{[1]a}.
 \tag{RRT.5}
\]
Then $w_a(F)\ge0$ for every polynomial, it vanishes off the stated
labels, and its sum over all monic degree-$N$ polynomials is $1$.
These assertions apply to the original \texttt{wtOf}; no
\texttt{admLq} hypothesis is needed for them.

For the proof, left multiplication by $b\in\mathbb C[G]$ has
diagonal entry $[1]b$ at every vector of the group basis.  Hence
\[
                 \operatorname{tr}L_b=|G|[1]b.
 \tag{RRT.6}
\]
Since $a$ and $E_F$ commute and are idempotent, so is $aE_F$.
Any idempotent linear map is the identity on its image and zero on
its kernel.  Its trace is its rank, a nonnegative integer.
Therefore $[1](aE_F)$ is real and nonnegative.  Its denominator
is positive by RRT1.  The vanishing assertion follows from
$E_F=0$; summing $aE_F$ gives $a$, proving normalization.
In fact
\[
 w_a(F)=\frac{\dim\operatorname{im}L_{aE_F}}
                     {\dim\operatorname{im}L_a}.
 \tag{RRT.7}
\]
The monic coefficient parametrization
$F_c=X^N+\sum_{i<N}c_iX^i$ is a bijection from
$\mathbb F_q^N$ to all monic degree-$N$ polynomials.  Thus
RRT2 proves exactly the finite probability normalization used by
\texttt{ConcNWeight.sum\_wtOf}.
\hfill$\square$
% END RRT2

% BEGIN RRT3
\paragraph{RRT3. The remaining character calculations have precise multiplicity targets.}
\label{claim:RRT3}
For an irreducible complex representation $\tau$ of $G$, let
$d_\tau=\dim\tau$, let $z_\tau$ be its primitive central
idempotent, and write $m_a(\tau)=\dim\operatorname{im}\tau(a)$.
The first identity below holds for any idempotent $a$; the
weight identity is used for $a=e_U,e_J$, whose denominators
are positive:
\[
 [1](az_\tau)=\frac{d_\tau m_a(\tau)}{|G|},\qquad
 w_a(F)=\frac{\sum_{\tau\in\mathcal S_F}d_\tau m_a(\tau)}
                                      {|G|[1]a},
 \tag{RRT.8}
\]
where $\mathcal S_F$ is the set of irreducibles on which $E_F$
acts as $1$.  To see the first equality directly, the coefficient
formula
$z_\tau(g)=d_\tau|G|^{-1}\operatorname{tr}\tau(g^{-1})$
gives $[1](az_\tau)=d_\tau|G|^{-1}\operatorname{tr}\tau(a)$.
The trace of the idempotent $\tau(a)$ is $m_a(\tau)$.
Maschke's decomposition, obtained by averaging a projection onto
any invariant subspace, and Schur's lemma show that each central
idempotent is the sum of the $z_\tau$ on which it acts as $1$.
This proves the second equality.

Consequently the actual lower comparison
\[
       (1-q^{-1})w_{e_U}(F)\le w_{e_J}(F)
 \tag{RRT.9}
\]
follows if for each irreducible in that same series
\[
       (q^{k_0}-q^{k_0-1})m_{e_U}(\tau)\le m_{e_J}(\tau).
 \tag{RRT.10}
\]
Indeed, multiply (RRT.10) by $d_\tau$, sum it, and use the
ratio $[1]e_J/[1]e_U=q^{k_0}$ from RRT1.
Likewise, for a squarefree label $F=\prod P$, the two character
facts
\[
 \mathcal S_F=\{\tau_F\},\qquad
 m_{e_U}(\tau_F)=1,\qquad
 d_{\tau_F}=\frac{\prod_{j=1}^N(q^j-1)}{\prod_{P\mid F}(q^{\deg P}-1)}
 \tag{RRT.11}
\]
give the exact squarefree law
\[
                w_{e_U}(F)=\prod_{P\mid F}(q^{\deg P}-1)^{-1}.
 \tag{RRT.12}
\]
Here $|G|=q^M\prod_{j=1}^N(q^j-1)$, by counting ordered
bases, and $[1]e_U=q^{-M}$.  Statements (RRT.10)--(RRT.11)
are explicit remaining character inputs, not conclusions of
idempotence.  The source proves them using the mirabolic trace
calculation, root exchange, and the squarefree Harish--Chandra
series theorem.  The squarefree \emph{upper} comparison for $e_J$
is not used in the probability proof below.
\hfill$\square$
% END RRT3

% BEGIN RRT4
\paragraph{RRT4. One-sided domination controls every product event.}
\label{claim:RRT4}
Let $P_v,Q_v$ be probability laws on finite sets, indexed by a
set of size $f$, and suppose $Q_v\ge(1-\delta)P_v$ pointwise,
where $0\le\delta\le1$.  For every event $B$ in their product,
\[
 \left|\bigotimes_v Q_v(B)-\bigotimes_v P_v(B)\right|
       \le1-(1-\delta)^f\le f\delta.
 \tag{RRT.13}
\]
When $\delta>0$, write
$Q_v=(1-\delta)P_v+\delta R_v$; the remainder is nonnegative
and sums to $1$.  Expanding the product gives
$\bigotimes Q_v=(1-\delta)^f\bigotimes P_v+
(1-(1-\delta)^f)R$ for a probability law $R$, except when the
second coefficient is zero, in which case the two laws already
agree.  Their event probabilities differ by at most that
coefficient.  The last inequality follows by induction from
$(1-\delta)^{f+1}\ge1-(f+1)\delta$.  The cases $\delta=0$ and
$f=0$ are immediate.

For the actual local laws, (RRT.9) and RRT2 apply with
$\delta=1/q$, $P_v=w_{e_U}$, and $Q_v$ equal to either
$w_{e_U}$ or $w_{e_J}$.  Thus the cost of transferring any bad
event to the actual sharp-level product law is at most $f/q$.
This proves the source's \texttt{ComparisonTransfer}, with a
slightly sharper first bound, from the lower comparison alone.
\hfill$\square$
% END RRT4

\subsection{The finite polynomial estimates}

For a monic squarefree polynomial $F$ with $F(0)\ne0$ put
\[
           w(F)=\prod_{P\mid F}(q^{\deg P}-1)^{-1}.
 \tag{RRT.14}
\]
The product runs over its distinct monic irreducible factors.
Let $\mathcal S_N$ denote these polynomials of degree $N$; put
$\mathcal S_0=\{1\}$.  These are finite sets, so all sums below
are ordinary finite sums.  They agree with the source's
\texttt{sqSet} and \texttt{sqWeight} under coefficient
parametrization.

% BEGIN RRT5
\paragraph{RRT5. Factor degrees are dominated by uniform permutation cycle degrees.}
\label{claim:RRT5}
For any property $\mathcal A$ of multisets of positive integers,
\[
 \sum_{\substack{F\in\mathcal S_N\\
               \mathcal A(\{\deg P:P\mid F\})}}w(F)
 \le\mathbb P_{\sigma\in S_N}
              \bigl(\mathcal A(\text{cycle lengths of }\sigma)\bigr).
 \tag{RRT.15}
\]
To prove this, let $I_d$ count the monic irreducibles of degree
$d$ with nonzero constant coefficient.  Their roots lie in
$\mathbb F_{q^d}^\times$, and each contributes $d$ distinct roots
disjoint from the roots of the other irreducibles.  Hence
$dI_d\le q^d-1$.  This root assertion follows either in the
quotient field $\mathbb F_q[X]/(P)$, which has $q^d$ elements,
or by its embedding in an algebraic closure; finite fields are
perfect and every nonzero element satisfies $x^{q^d-1}=1$.

For prescribed counts $a_d\ge0$ with $\sum_d da_d=N$, unique
factorization gives exactly $\prod_d\binom{I_d}{a_d}$ such
polynomials.  Their total weight is
\[
 \prod_d\frac{\binom{I_d}{a_d}}{(q^d-1)^{a_d}}
       \le\prod_d\frac1{d^{a_d}a_d!}.
 \tag{RRT.16}
\]
The right side is the probability of that cycle type for a
uniform permutation: arranging the $N$ symbols gives $N!$
orders, and one divides by the $d$ rotations of each $d$-cycle
and by the $a_d!$ permutations of equally long cycles.
Summing (RRT.16) over the allowed types proves (RRT.15).
In particular $\sum_{F\in\mathcal S_N}w(F)\le1$.

There is also the useful exact divisor consequence.  If $H$ is
monic squarefree, $H(0)\ne0$, and $\deg H\le N$, then
\[
                 \sum_{\substack{F\in\mathcal S_N\\H\mid F}}w(F)
                         \le w(H).
 \tag{RRT.17}
\]
Indeed $F=HR$ with $R\in\mathcal S_{N-\deg H}$ and
$\gcd(H,R)=1$, and $w(F)=w(H)w(R)$.  Dropping the coprimality
condition and using the preceding mass bound proves (RRT.17).
Thus \texttt{TypeWeightExact}, \texttt{TypeSetWeight},
\texttt{SqMassLeOne}, and \texttt{DivisorBound} all have direct
human proofs here.
\hfill$\square$
% END RRT5

% BEGIN RRT6
\paragraph{RRT6. The squarefree exceptional mass.}
\label{claim:RRT6}
For $N\ge1$,
\[
                  \sum_{F\in\mathcal S_N}w(F)\ge1-2/q.
 \tag{RRT.18}
\]
There are $q^N$ monic degree-$N$ polynomials, and at most
$q^{N-1}$ have zero constant coefficient.  For $N\ge2$ the
number of monic squarefree degree-$N$ polynomials is
$q^N-q^{N-1}$.  Here is an elementary count: every monic
polynomial uniquely has the form $AB^2$ with $A$ squarefree,
by splitting each irreducible-factor exponent into its parity
and twice an integer.  If $S(t)$ is the formal generating series
for monic squarefree polynomials, coefficientwise finite
counting gives
\[
 \frac1{1-qt}=S(t)\frac1{1-qt^2},\qquad
                  S(t)=\frac{1-qt^2}{1-qt}.
\]
Its coefficients are $1,q,q^2-q,q^3-q^2,\ldots$.
Thus $|\mathcal S_N|\ge q^N-2q^{N-1}$; for $N=1$,
$|\mathcal S_1|=q-1$ gives the same lower bound.
Each $F\in\mathcal S_N$ has $w(F)\ge q^{-N}$, because the
sum of its factor degrees is $N$.  This proves (RRT.18).

In particular, any probability law $P$ on monic degree-$N$
polynomials that equals $w$ on $\mathcal S_N$ has
$P(\mathcal S_N^c)\le2/q$.  In $f$ independent coordinates,
the probability that some coordinate is outside $\mathcal S_N$
is at most $2f/q$.  Combining this with RRT4 accounts exactly
for the term $3f/q$ in the original $W_g$.
\hfill$\square$
% END RRT6

% BEGIN RRT7
\paragraph{RRT7. A permutation has an invariant proper size with probability at most $7/8$.}
\label{claim:RRT7}
Fix $1\le a<N$.  For a permutation $\sigma$, let $Z(\sigma)$
be its number of invariant $a$-element subsets.  Let
$Y(\sigma)$ count those invariant subsets containing a fixed
point and having a fixed point in their complement.  Uniformly
over $S_N$,
\[
                \mathbb EZ=1,\qquad \mathbb EY\ge1/4.
 \tag{RRT.19}
\]
For the first identity, any particular $a$-set is preserved by
$a!(N-a)!$ permutations.  For the second, conditional on
preserving that set, the two restrictions are independent uniform
permutations.  A uniform permutation of any nonempty finite set
has a fixed point with probability at least $1/2$.  For sets of
size one this is immediate; otherwise inclusion--exclusion gives
the derangement probability $\sum_{j=0}^n(-1)^j/j!\le1/2$:
after the first three terms, consecutive negative--positive
pairs have nonpositive sum, and a possible remaining term is
negative.  Summing over all $a$-sets proves (RRT.19).

Pointwise we have
\[
                 2\mathbf1_{Z>0}+Y\le2Z.
 \tag{RRT.20}
\]
If $Y=0$, this follows from $Z\ge1$ on the indicated event.
If $Y>0$, exchange one fixed point in a counted set with a fixed
point outside it.  This gives a different invariant $a$-set,
so $Z\ge2$; since $Y\le Z$, (RRT.20) again follows.
Taking expectations yields
$2\mathbb P(Z>0)+1/4\le2$, hence
\[
                       \mathbb P(Z>0)\le7/8.
 \tag{RRT.21}
\]
This proves the full original \texttt{LemmaP}, without a
large-$N$ restriction.
\hfill$\square$
% END RRT7

% BEGIN RRT8
\paragraph{RRT8. No common proper divisor degree, with the source constant.}
\label{claim:RRT8}
For $1\le a<N$, a squarefree $F$ has a divisor of degree $a$
if and only if some subset of its irreducible-factor degrees
sums to $a$.  A permutation has an invariant $a$-set if and
only if a subset of its cycle lengths sums to $a$.
RRT5 and RRT7 therefore give
\[
           \sum_{\substack{F\in\mathcal S_N\\
                      \exists H\mid F:\deg H=a}}w(F)\le7/8.
 \tag{RRT.22}
\]
This is the original \texttt{LocalEstimateIV}.  For any
probability law $P$ agreeing with $w$ on $\mathcal S_N$, the
probability that all $f$ independent coordinates are squarefree
and have a divisor of this same degree is at most $(7/8)^f$.
There are $\lfloor N/2\rfloor$ possible degrees in the
source's common-degree obstruction.  The union bound gives
\[
   \mathbb P(\text{squarefree common-degree obstruction})
                          \le\frac N2(7/8)^f.
 \tag{RRT.23}
\]
No assertion about the nonsquarefree part of $P$ is needed for
this event; that part was handled separately in RRT6.
\hfill$\square$
% END RRT8

% BEGIN RRT9
\paragraph{RRT9. An elementary eventual long-prime reciprocal bound.}
\label{claim:RRT9}
There is an absolute $N_0\ge100$ such that for every integer
$N\ge N_0$,
\[
       \sum_{\substack{p\text{ prime}\\N/2<p,\ p+3\le N}}
                           \frac1p\ge\frac1{5\log N}.
 \tag{RRT.24}
\]
Here only an eventual version is needed; the source proves the
stronger statement starting at $N=100$ by additional finite
checks.  The proof below avoids those checks and keeps the same
constant $5$.

First, the primorial bound $\prod_{p\le m}p\le4^m$ follows by
strong induction.  For $m=2n$, the product over $n<p\le2n$
divides $\binom{2n}{n}\le4^n$, and one uses the inductive bound
at $n$.  For $m=2n+1$, the product over
$n+1<p\le2n+1$ divides $\binom{2n+1}{n}\le4^n$; the last
inequality follows because its two equal middle terms already
contribute twice that coefficient to the sum $2^{2n+1}$.
Use the inductive bound at $n+1$.  The initial cases are direct.

Let $b_n$ count the primes in $(n,2n]$.  Legendre's factorial
formula gives
\[
 \frac{4^n}{2n+1}\le\binom{2n}{n}
                 \le (2n)^{\sqrt{2n}}4^{2n/3}(2n)^{b_n}
 \tag{RRT.25}
\]
for all sufficiently large $n$.  To check the upper bound,
the exponent of $p$ in the middle binomial coefficient is a sum
of terms
$\lfloor2n/p^j\rfloor-2\lfloor n/p^j\rfloor\in\{0,1\}$.
Thus the full power contributed by each $p\le\sqrt{2n}$ is at
most $2n$, and there are at most $\sqrt{2n}$ such primes.
Above $\sqrt{2n}$, each exponent is at most one.  The primes up
to $2n/3$ are bounded by the primorial estimate, the primes in
$(2n/3,n]$ have exponent zero, and the remaining $b_n$ primes
contribute at most $(2n)^{b_n}$.  The lower bound in (RRT.25)
comes from the largest of the $2n+1$ binomial coefficients.

Taking logarithms gives
\[
 b_n\ge\frac{(n/3)\log4-\log(2n+1)
                           -\sqrt{2n}\log(2n)}{\log(2n)}.
 \tag{RRT.26}
\]
Set $n=\lfloor N/2\rfloor$.  Every prime in $(n,2n]$ exceeds
$N/2$, and excluding the last three possible integers leaves
at least $b_n-3$ primes satisfying $p+3\le N$.
Each contributes at least $1/N$.  The lower bound for
$(b_n-3)\log N/N$ furnished by (RRT.26) tends to
$\log4/6>1/5$.  For instance
$\log t\ge2(t-1)/(t+1)$ for $t\ge1$ follows by differentiating
the difference; applying its strict inequality at $t=2$ gives
$\log4>4/3$.  Since $\log N/\sqrt N\to0$, one absolute
threshold makes the bound at least $1/5$, proving (RRT.24).
All estimates are unconditional elementary integer counts.
\hfill$\square$
% END RRT9

% BEGIN RRT10
\paragraph{RRT10. The long-prime-degree witness appears exponentially often.}
\label{claim:RRT10}
For $N\ge N_0$ and a probability law $P$ agreeing with $w$ on
$\mathcal S_N$,
\[
 P\{F\in\mathcal S_N:\text{no irreducible factor has prime degree }p,
                         \ N/2<p,\ p+3\le N\}
                   \le1-\frac1{5\log N}.
 \tag{RRT.27}
\]
Indeed, the expected number of $p$-cycles in a uniform
permutation is $1/p$: choose the $p$ symbols, arrange them into
one of $(p-1)!$ cycles, and permute the remaining symbols.
Cycles of lengths exceeding $N/2$ cannot coexist.  Hence the
probability of at least one of the specified prime cycles is
exactly the reciprocal sum in (RRT.24).  Apply RRT5 to the
\emph{failure} property, and then RRT9.  This proves the
eventual version of the original \texttt{LocalEstimateV}.
Independence across the $f$ places and $1-t\le e^{-t}$ give
\[
 P^{\otimes f}\{\text{every coordinate is squarefree and lacks
                          this witness}\}
                       \le\exp\{-f/(5\log N)\}.
 \tag{RRT.28}
\]
The event includes squarefreeness; other failures were already
assigned their separate cost in RRT6.
\hfill$\square$
% END RRT10

\subsection{The exact remaining probability inputs and the original consumer}

For a nonzero element $z$ of an algebraic closure of a finite
field, $z_{\ell'}$ denotes the unique prime-to-$\ell$ part of
its finite order.  If $\operatorname{ord}(z)=\ell^a m$ with
$(\ell,m)=1$, it is the power of $z$ whose exponent is $0$
modulo $\ell^a$ and $1$ modulo $m$.  This characterizes both
source formulas for \texttt{primeToPart}: their exponents have
these same two residues.  All roots considered below are
nonzero, and lie in a finite field, so this comparison applies.
Put
\[
 \mathcal B(q,N)=\{q\}\cup
 \{\ell\text{ prime}:\exists\,1\le d\le N,
                           \ q^d<\ell^{4v_\ell(q^d-1)}\}.
 \tag{RRT.29}
\]
For prime $q$ this is exactly either source version of
\texttt{BB}.  The formally finite range loses nothing: a prime
in the displayed second set divides some $q^d-1$, hence is at
most $q^N$.

Here are the two exact collision inequalities on this same
polynomial space.  The first remains an input.  The second
is proved completely in RRT14--RRT16 below:
\begin{align}
 \sum_{\substack{F\in\mathcal S_N\\
   \exists\ell\notin\mathcal B(q,N):\\
   z\mapsto z_{\ell'}\text{ is not injective on roots of }F}}
 w(F)&\le\frac{23\log q}{\sqrt q},                 \tag{RRT.30}\\
 \sum_{\substack{F\in\mathcal S_N\\
   \exists\ell\ne q\text{ prime},\ z\text{ root of }F:\\
                         z_{\ell'}=c_{\ell'}}}w(F)
 &\le A_q:=\frac1{4\min(q_1,d_0)}
                  +\frac{1+\log N}{q-1}\quad(c\in\mathbb F_q^\times).
                                                        \tag{RRT.31}
\end{align}
In (RRT.30) the quantified $\ell$ is prime.  The exact source
guards for these two inputs include $q\ge2^{12}$,
$q_1,d_0$ distinct odd primes, $\max(q_1,d_0)\ge16$, and
$16,q_1,d_0\mid q-1$.  The two statements are respectively
\texttt{LocalEstimateII} and \texttt{LocalEstimateIII}.
RRT5's divisor bound is one of their now proved ingredients.

% BEGIN RRT11
\paragraph{RRT11. Assemble the original $W_g$ on a reduced local contract.}
\label{claim:RRT11}
Let $\iota$ have $f$ elements, let the twist set have $m$
elements, and prescribe $\chi_{v,j}\in\mathbb F_q^\times$.
Assume $N\ge N_0$, the guards just stated, the original
series laws (RRT.9), (RRT.12), and the root-injectivity
estimate (RRT.30).  The second collision estimate (RRT.31)
is supplied by RRT16.  Choose each actual place law $Q_v$ to be
$w_{e_J}$ or $w_{e_U}$.  Then the probability of a tuple failing
the source's \texttt{goodAt} condition is at most
\[
 \frac{3f}{q}+\frac{23f\log q}{\sqrt q}
       +2fmA_q+\frac N2(7/8)^f
                         +\exp\{-f/(5\log N)\}.
 \tag{RRT.32}
\]

To prove this, first use the common base law $P=w_{e_U}$.
RRT2 proves that it is a probability on the actual labels;
(RRT.12) identifies its squarefree weights.  A tuple satisfies
the six clauses of \texttt{GoodAt} precisely when its
polynomials are valid labels, all are squarefree, their roots
remain distinct after every allowed prime-to-$\ell$ projection,
no root collides under that projection with any
$\chi_{v,j}$ or $\chi_{v,j}^{-1}$, every degree
$1\le a\le N/2$ is absent as a divisor degree at some place,
and some place supplies the long prime-degree factor of RRT10.
The prime-to-part and exceptional-set identifications above
identify this with the conjunction of the two original
\texttt{IsGood} predicates defining \texttt{goodAt}.

The failure set is contained in the union of: a nonsquarefree
coordinate; a root-injectivity collision at a coordinate; one
of the two twist collisions at one of the $fm$ pairs; a common
proper divisor degree across all coordinates; and absence of
the long-prime-degree witness at every coordinate.  RRT6,
(RRT.30), (RRT.31), RRT8, and RRT10 bound these probabilities
by the respective terms in (RRT.32), with $2f/q$ in place of
$3f/q$.  Independence is needed only for the last two events;
the other three use the union bound.  RRT4 adds at most $f/q$
when passing to $\bigotimes Q_v$, proving (RRT.32).

For the actual datum take $q=g_k$, $N=R.r$, $f=R.f$,
$m=R.m$, $(q_1,d_0)=(\mathrm{host}.q,\mathrm{host}.d)$, and
the original twist table.  The cardinality field
$|\mathrm{PlG}(k)|=R.f$ makes (RRT.32) the literal $W_g$.
RRT2 supplies \emph{all} local nonnegativity assertions in
\texttt{WeightAllRows}.  Any finite set of bad tuples has
weight at most the full bad event, so (RRT.32) supplies its
second assertion as well.  Polynomial tuples that are not
valid labels contribute zero; coefficient parametrization is
injective on the others.  Repeated global row labels do not
affect this local argument.

The field \texttt{AllJoinA.level} gives $Q_v=w_{e_J}$ at
every place, with one nontrivial additive character per
auxiliary prime.  Its bundle field is not used in this local
weight argument.  Only the lower comparison and squarefree
$e_U$ formula are required from \texttt{admLq}; its
squarefree upper bound on $e_J$ is unnecessary for this
consumer.  Finally $R.f\le R.k\le R.r$ and
$x^2\le R.f$ imply $R.r\ge N_0$ after the single absolute
threshold $x\ge\sqrt{N_0}$.  Thus RRT9's eventual version
suffices for the original eventual \texttt{LinkIRR}
application, with exactly the same $W_g$ and all choice orders
of GQT5.  Here the rank inequalities come from
\texttt{Params.lean}: $R.f=2R.\mathrm{ord}$,
$2R.\mathrm{ord}\le R.k$, $R.r=R.k+R.m-1$, and $R.m\ge2$.
They are explicitly used again in \texttt{sizesRec\_at} of
\texttt{LinkIRROfJoinEx.lean}.  The abstract numerical
\texttt{SizesAt} conditions alone do not assert $N\to\infty$;
the displayed actual-rank comparison is essential here.
\hfill$\square$
% END RRT11

\subsection{Two algebraic parts of the actual row construction}

% BEGIN RRT12
\paragraph{RRT12. Exact finite tensor cylinders.}
\label{claim:RRT12}
Let the finite local groups be $G_v$, with central orthogonal
idempotents $E_{v,t}$ summing to $1$, and nonzero idempotents
$a_v$ having positive identity coefficients.  In
$\mathbb C[\prod_vG_v]\simeq\bigotimes_v\mathbb C[G_v]$ put
\[
 a=\bigotimes_va_v,\qquad
 E_S=\sum_{(t_v)\in S}\bigotimes_vE_{v,t_v}.
 \tag{RRT.33}
\]
Then $E_S$ is central and idempotent,
$E_SE_T=E_{S\cap T}$, and $E_{\rm all}=1$.
For an event determined by a subcollection $V_0$ of places,
\[
 \frac{[1](aE_S)}{[1]a}
          =\sum_{(t_v)_{v\in V_0}\in S}
                   \prod_{v\in V_0}w_{a_v}(t_v).
 \tag{RRT.34}
\]
Indeed products of different tuple idempotents vanish and
their sums expand coordinate by coordinate.  Identity
coefficients multiply under the displayed tensor product.
The complementary coordinates sum to $a_v$, so their factors
cancel in the ratio.  RRT2 makes the remaining factors real
and nonnegative.  These steps prove every assertion directly.

For the two auxiliary-prime collections, (RRT.33) is exactly
the finite \texttt{AllRowsOf.ES} coefficient formula.  Its
universal event and cylinder value are therefore proved on
the original finite groups.  Passing to the original adelic
function $e_T$, including its Haar factor, still requires the
actual reduction/inflation and convolution comparison
\texttt{EtOne}; (RRT.34) does not choose or rescale that Haar
measure.  It supplies the finite algebraic identity inside
that comparison.
\hfill$\square$
% END RRT12

% BEGIN RRT13
\paragraph{RRT13. Construct labels and the retained injection without distinct-label assumptions.}
\label{claim:RRT13}
Let $L$ be the finite set of valid tuple labels, with good
subset $L^+$.  Suppose exact nonnegative integer multiplicities
$d_t$ satisfy $\sum_{t\in L}d_t=D_{\rm all}$ and that the
given retained rows $\sigma\in\{1,\ldots,D_{\rm rows}\}$ have
their original good labels $\lambda_\sigma\in L^+$, with
\[
       |\{\sigma:\lambda_\sigma=t\}|=d_t\qquad(t\in L^+).
 \tag{RRT.35}
\]
Then there is a labeling of $D_{\rm all}$ rows and an
injection of the original retained rows preserving every
label, whose image is exactly the good rows, and with exactly
$\sum_{t\in S}d_t$ rows in each set of labels $S$.

Construct the finite set
\[
 T=\{1,\ldots,D_{\rm rows}\}\ \sqcup\!
            \coprod_{t\in L\setminus L^+}\{1,\ldots,d_t\}.
 \tag{RRT.36}
\]
Keep $\lambda_\sigma$ on the first summand and assign label
$t$ to the $t$-summand on the right.  Summing (RRT.35) over
$L^+$ shows $|T|=D_{\rm all}$.  Choose any enumeration of
$T$ by the full row index set and transport this labeling.
The first inclusion gives the required injection.  Its image
is exactly the good part because every other summand has a
bad label.  Fiber sizes are (RRT.35) on good labels and the
specified $d_t$ on bad labels.  Summing those fibers over $S$
gives the assertion.  This construction includes zero
multiplicities and repeated labels, and does not require
choosing distinct polynomials for distinct rows.

In the original construction take
$d_t=\mathrm{dimX}(\{t\})$.  The hypotheses
\texttt{RowsAdd}, \texttt{RowsGood}, \texttt{EtUniv}, and
the original \texttt{AutTie} give the exact additivity,
good-fiber equation, and total rank needed above.  Use the
original retained labels
$\lambda_\sigma(k)=D.\mathrm{nd}.\mathrm{lab}(k,
D.\mathrm{what}_2.\mathrm{ft}(\sigma))$.
RRT13 then constructs \texttt{RowsLab}, with all its fields,
and hence the original \texttt{LabelCountGoodInjective} once
the same \texttt{allRowsOf} object is installed.  The
representation-theoretic producers of \texttt{RowsAdd} and
\texttt{RowsGood} are separate requirements, not consequences
of finite enumeration.
\hfill$\square$
% END RRT13

\subsection{The complete root-to-twist collision estimate}

% BEGIN RRT14
\paragraph{RRT14. Two prime-power sums, with the original constants.}
\label{claim:RRT14}
For a positive integer $n$ write
$S(n)=\sum_{\ell\mid n,\ \ell\text{ prime}}\ell^{v_\ell(n)}$.
If $q\ge2$, $4\mid q-1$, and $\delta\ge1$, then
\[
 S(q^\delta-1)\le\delta(q-1)+\frac{q^\delta-1}{q-1}.
 \tag{RRT.37}
\]
If moreover $p_1,p_2$ are distinct odd primes dividing $q-1$,
$16\mid q-1$, and $\max(p_1,p_2)\ge16$, then
\[
               \frac{S(q-1)}{q-1}
                         \le\frac1{4\min(p_1,p_2)}.
 \tag{RRT.38}
\]
These are exactly \texttt{PrimePowerOrderCount} and
\texttt{PrimePowerOrderShare}.

For (RRT.37), first observe that for any finite collection of
integers at least $2$, its sum is at most its product; this
includes the empty collection, whose sum is $0$ and product
is $1$.  For a nonempty collection, induction uses
$u+v\le uv$ for $u,v\ge2$.  Put $A=q-1$ and
$B=(q^\delta-1)/A$.  For a prime $\ell\mid A$,
\[
 v_\ell(q^\delta-1)=v_\ell(A)+v_\ell(\delta).
 \tag{RRT.39}
\]
Here is the needed elementary valuation proof.  If
$u=1+\ell^h b$ with $\ell\nmid b$, and either $\ell$ is odd
and $h\ge1$ or $\ell=2$ and $h\ge2$, expansion of
$u^\ell-1$ has first term $\ell^{h+1}b$ and every other
term has valuation strictly greater than $h+1$.  Hence
raising to the power $\ell$ raises the valuation by exactly
one.  Raising to an integer coprime to $\ell$ keeps valuation
$h$, again by the first binomial term.  Splitting $\delta$
into an $\ell$-power times a coprime integer proves (RRT.39).
For $\ell=2$ the hypothesis $4\mid A$ gives $h\ge2$.

The prime-power factors of $q^\delta-1$ at primes dividing
$A$ therefore have product dividing $\delta A$.  Their sum
is at most that product, hence at most $\delta A$.  The other
prime-power factors have product dividing $B$, since these
primes do not divide $A$; their sum is at most $B$.
Adding proves (RRT.37), including empty parts.

For (RRT.38) set $m=\min(p_1,p_2)$ and isolate the
prime-power factors $X$ at $2$, $Y$ at the smaller selected
prime, and $Z$ at the larger selected prime.  Then
$X\ge16$, $Y\ge m$, $Z\ge16$, and $Z\ge m$.
Let $P$ be the product of all remaining prime-power factors,
and $T$ their sum.  Thus $P\ge1$, $T\le P$,
$q-1=XYZP$, and $S(q-1)=X+Y+Z+T$.
Each of
\[
 \frac{X}{XYZP},\qquad \frac{Y}{XYZP},\qquad
 \frac{Z}{XYZP},\qquad \frac{T}{XYZP}
\]
is at most $1/(16m)$: use respectively $YZ\ge16m$,
$XZ\ge16m$, $XY\ge16m$, and $T\le P$ together with
$XYZ\ge16m$.  Their sum proves (RRT.38).
\hfill$\square$
% END RRT14

% BEGIN RRT15
\paragraph{RRT15. Count the actual irreducible collision factors.}
\label{claim:RRT15}
Fix $c\in\mathbb F_q^\times$ and $\delta\ge1$, with $q\ge3$.
Let $\mathcal B_\delta(c)$ be the monic irreducible polynomials
of degree $\delta$ that divide at least one polynomial
\[
 X^{\ell^{v_\ell(q^\delta-1)}}
             -c^{\ell^{v_\ell(q^\delta-1)}}
       \quad(\ell\text{ prime},\ \ell\mid q^\delta-1).
 \tag{RRT.40}
\]
Write $b_\delta=|\mathcal B_\delta(c)|$.  Then
\[
                  \delta b_\delta\le S(q^\delta-1).
 \tag{RRT.41}
\]
Every monic irreducible degree-$\delta$ polynomial with a
nonzero root $z$ satisfying $z_{\ell'}=c_{\ell'}$ for some
prime $\ell$ belongs to $\mathcal B_\delta(c)$.

Indeed, $z$ and $c$ lie in $\mathbb F_{q^\delta}^\times$.
The canonical prime-to-$\ell$ map is a homomorphism on this
finite cyclic group, so the quotient $z/c$ has $\ell$-power
order.  Its order also divides $q^\delta-1$.  If that order
is greater than one, it divides
$\ell^{v_\ell(q^\delta-1)}$.  If it is one, use any prime
divisor of $q^\delta-1$, which exists since $q\ge3$.
Thus $z$ is a root of one polynomial in (RRT.40), and its
monic minimal polynomial divides that polynomial.

For each fixed $\ell$, the distinct monic irreducible factors
of degree $\delta$ in (RRT.40) are coprime.  Their product
divides its nonzero polynomial of degree
$\ell^{v_\ell(q^\delta-1)}$.  Consequently their number
times $\delta$ is at most this degree.  Summing these bounds
over $\ell$ proves (RRT.41).  Counting a factor more than
once only enlarges the bound.  Every polynomial in
$\mathcal B_\delta(c)$ has nonzero constant coefficient,
because the polynomials in (RRT.40) do.
\hfill$\square$
% END RRT15

% BEGIN RRT16
\paragraph{RRT16. The full original LocalEstimateIII.}
\label{claim:RRT16}
Assume $N\ge1$, $p_1,p_2$ are distinct odd primes,
$16,p_1,p_2\mid q-1$, and $\max(p_1,p_2)\ge16$.
Then for every $c\in\mathbb F_q^\times$,
\[
 \sum_{\substack{F\in\mathcal S_N\\
        \exists\ell\ne q\text{ prime},\ z\text{ root of }F:\\
                                      z_{\ell'}=c_{\ell'}}}w(F)
        \le\frac1{4\min(p_1,p_2)}+\frac{1+\log N}{q-1}.
 \tag{RRT.42}
\]
In fact the proof also allows the prime $\ell=q$ whenever
$q$ is prime.  The displayed weaker form is the exact
original \texttt{LocalEstimateIII}.

If $F$ occurs in this sum, the monic minimal polynomial of
its indicated root is a factor of $F$, has a degree
$1\le\delta\le N$, and belongs to $\mathcal B_\delta(c)$
by RRT15.  RRT5's divisor bound, followed by a union bound
over these factors, therefore bounds the left side by
\[
                  \sum_{\delta=1}^N
                                \frac{b_\delta}{q^\delta-1}.
 \tag{RRT.43}
\]
For $\delta=1$, RRT15 and (RRT.38) bound this summand by
$1/(4\min(p_1,p_2))$.  For $\delta\ge2$, RRT15 and
(RRT.37) give
\[
 \frac{b_\delta}{q^\delta-1}
       \le\frac{q-1}{q^\delta-1}+\frac1{\delta(q-1)}
       \le q^{-(\delta-1)}+\frac1{\delta(q-1)}.
 \tag{RRT.44}
\]
The last comparison follows from
$(q-1)q^{\delta-1}\le q^\delta-1$.
Finally
\[
 \sum_{\delta=2}^Nq^{-(\delta-1)}\le\frac1{q-1},
 \qquad \sum_{\delta=2}^N\frac1\delta\le\log N.
\]
The first is the finite geometric sum; the second follows
by bounding $1/\delta$ by
$\int_{\delta-1}^{\delta}dt/t$.  They also hold for $N=1$
with empty sums.  Substitution in (RRT.43) proves (RRT.42).
There is no character-theoretic hypothesis in this proof:
the weight is the explicit squarefree product (RRT.14).
Installing $p_1=q_1,p_2=d_0$ proves (RRT.31) and supplies
the entire $2fmA_q$ term in RRT11 once the base probability
law is identified by (RRT.12).
\hfill$\square$
% END RRT16

\subsection{What has advanced, and the remaining frontier}

The nonnegativity and normalization part of the actual local
weights is now proved without Green.  The source's
\texttt{ComparisonTransfer}, \texttt{TypeSetWeight},
\texttt{SqMassLower}, \texttt{SqMassLeOne},
\texttt{DivisorBound}, \texttt{LemmaP}, and
\texttt{LocalEstimateIII} and \texttt{LocalEstimateIV} have
complete finite proofs here.
The eventual \texttt{LocalEstimateV} has the original constant
and is enough for the selected family.  This supplies four of
the five displayed terms of $W_g$, including their actual
probabilistic meaning, and the exact finite tensor and label
constructions required upstream of GQT5.

Three local quantitative inputs remain: the original
squarefree $e_U$ character law, the $e_J$ lower comparison,
and the root-injectivity collision inequality (RRT.30).
For the first two, (RRT.10)--(RRT.11) are explicit finite
character targets.  The source's normalized-label bridge
permutes irreducible polynomial labels by degree before
identifying the original series projector with Green's
parametrization; it must be written on the same labels, not
assumed from a name.  Its exact entry point is
\texttt{AdmissibleLabelWeightsOfGreenCuspidalParametrisa}.
For the remaining collision input the next proofs are the
same-factor and cross-factor root-fiber counts and their
weighted divisor sums.  They can now use RRT5 directly.

The global frontier still includes the original adelic
\texttt{HidaSetting}, reduction/inflation and \texttt{EtOne},
the actual \texttt{RowsAdd}/\texttt{RowsGood} identifications,
and the cylinder counts.  RRT13 is the full finite enumeration
after those identifications; it does not construct an
automorphic spectrum.  GQT1--GQT5 continue from the resulting
original row object and the three sufficient counting tests.
No new proof of RH or GRH, and no completion of
\texttt{GLOBAL\_RETAINED\_DEFECT}, is claimed by this packet.
